% Gerada por scripts/importar_equivalencia_java.py. Não editar à mão.
\begin{table}[htbp]
\centering
\footnotesize
\caption{Concordância entre o protótipo Java e a implementação Python do Comparador de Registros em 500 pares sintéticos construídos para exercitar casos de fronteira; proporção de pares com subescore idêntico (tolerância de $10^{-9}$) e maior diferença absoluta}
\label{tab:equivalencia-java}
\begin{tabular}{lrr}
\toprule
Subescore & Concordância (\%) & Diferença máxima \\
\midrule
Paciente: primeiro fragmento & 91,6 & 1,0 \\
Paciente: último fragmento & 93,2 & 1,0 \\
Paciente: fragmentos iguais & 84,8 & 1,0 \\
Paciente: fragmentos raros & 56,8 & 1,0 \\
Paciente: fragmentos comuns & 39,2 & 1,0 \\
Paciente: Soundex & 88,8 & 0,8 \\
Paciente: abreviações & 99,2 & 0,25 \\
Mãe: primeiro fragmento & 92,8 & 1,0 \\
Mãe: último fragmento & 93,2 & 1,0 \\
Mãe: fragmentos iguais & 86,4 & 1,0 \\
Mãe: fragmentos raros & 48,0 & 2,0 \\
Mãe: fragmentos comuns & 70,8 & 2,0 \\
Mãe: Soundex & 89,6 & 0,8 \\
Mãe: abreviações & 99,2 & 0,25 \\
\midrule
Data: igualdade & 100,0 & 0,0 \\
Data: um dígito & 100,0 & 0,0 \\
Data: inversão do dia & 82,4 & 1,0 \\
Data: inversão do mês & 76,8 & 1,0 \\
Data: inversão do ano & 75,6 & 0,8 \\
\midrule
Nota final & 11,2 & 4,0 \\
\bottomrule
\end{tabular}
\par\smallskip
\begin{minipage}{\textwidth}
\footnotesize Nota: o lote não estima a concordância em dados reais. As tabelas de frequência são sintéticas e idênticas nos dois lados. Fonte: \path{tests/equivalencia/resultado_agregado.csv} do repositório do CdR, commit \texttt{0d5b411}, PR \#25 \cite{Jardim2024comparador}. Importação e geração: \path{scripts/importar_equivalencia_java.py} desta tese.
\end{minipage}
\end{table}
